% TOP_PROBLEM: {"id": 2100603, "title": "Open Problems in Integrable Systems — Poisson geometry and action-angle variables", "category_id": 11, "category": {"id": 11, "name": "dynamical_systems", "display_name": "Dynamical Systems", "description": "Problems about long-term behavior of deterministic systems, Hamiltonian dynamics, and periodic orbits.", "slug": "dynamical-systems", "order_index": 11, "created_at": "2026-07-31T15:26:25.671Z"}, "status": "partially_solved"}
% FIRST_POSTED: null
% ENTRY_KIND: statement_only
% Imported from: batch09.tex; UnsolvedMath ID 2100603
\documentclass[11pt]{article}
\usepackage[margin=25mm]{geometry}
\usepackage{fontspec}
\setmainfont{Latin Modern Roman}
\usepackage{amsmath,amssymb,mathtools}
\usepackage{unicode-math}
\setmathfont{Latin Modern Math}
\usepackage{xurl,hyperref}
\hypersetup{colorlinks=true,urlcolor=blue}
\setcounter{secnumdepth}{0}
\setlength{\parindent}{0pt}
\setlength{\parskip}{5pt}
\setlength{\emergencystretch}{3em}
\newcommand{\sref}[2]{\hyperlink{source:#1:#2}{[S#2]}}
\newcommand{\cataloguescope}[1]{\paragraph{Recorded target.}#1}
\begin{document}
% UnsolvedMath ID 2100603; source catalogue code AMR-020-0603
\setcounter{section}{2100602}
\section{Open Problems in Integrable Systems — Poisson geometry and action-angle variables}\label{2100603}
{\small\sffamily UnsolvedMath ID 2100603 · Dynamical Systems}
\subsection{Definitions and mathematical statement}

\paragraph{Shared notation.}$\mathbb N=\{1,2,\ldots\}$, and $\mathbb Z,\mathbb Q,\mathbb R,\mathbb C$ denote the integers, rationals, reals and complex numbers. Any other conventions are specified locally.


Let $(M^N,\Pi)$ be a regular Poisson manifold of rank $2r$. Put $s=N-r$. A regular Liouville integrable system is a map $F=(f_1,\ldots,f_s):M\to\mathbb R^s$ with independent differentials and $\{f_i,f_j\}=0$, whose Hamiltonian leaves are compact $r$-tori. Locally the action-angle theorem gives coordinates $(\theta_1,\ldots,\theta_r,\sigma_1,\ldots,\sigma_s)$ with $\theta_i\in\mathbb R/\mathbb Z$ and
\[\Pi=\sum_{i=1}^r\frac{\partial}{\partial\theta_i}\wedge\frac{\partial}{\partial\sigma_i};\]
the remaining $\sigma_{r+1},\ldots,\sigma_s$ are Casimirs. Global action-angle coordinates mean a global product-coordinate representation $M\simeq\mathbb T^r\times B$, $B\subseteq\mathbb R^s$ open, with the same Poisson formula and with the level-set foliation of $F$ given by projection to $B$.
\paragraph{Statement recorded in the supplied source.}
Determine the obstructions to patching local action-angle coordinates into global action-angle coordinates for regular Poisson integrable systems. Characterize when those obstructions vanish, taking account of the foliation by affine leaves on the moment-map image and the global transverse Casimir coordinates. \sref{2100603}{2}
\subsection{Short English statement}
What prevents local Poisson action-angle charts from combining into one global system of angle, action and transverse coordinates?
\subsection{Sources}
\begin{description}
\item[\hypertarget{source:2100603:1}{S1}] ULAM.AI and the UnsolvedMath Contributors, UnsolvedMath: A Curated Collection of Open Mathematics Problems, record 2100603 (AMR-020-0603). CC BY 4.0. \url{https://huggingface.co/datasets/ulamai/UnsolvedMath}.
\item[\hypertarget{source:2100603:2}{S2}] Alexey Bolsinov, Vladimir S. Matveev, Eva Miranda and Serge Tabachnikov, Open Problems, Questions, and Challenges in Finite-Dimensional Integrable Systems (2018), Problem 6.5 and its complete preceding definitions. Original question and full discussion. \url{https://arxiv.org/abs/1804.03737}.

\end{description}
\label{end:2100603}
\end{document}
